\documentclass[11pt,a4paper]{article}
\usepackage[margin=2.5cm]{geometry}
\usepackage{times}
\usepackage{latexsym}
\usepackage{amsmath,amssymb,amsfonts}
\usepackage{booktabs}
\usepackage{multirow}
\usepackage{graphicx}
\usepackage{xcolor}
\usepackage{tabularx}
\usepackage{array}
\usepackage{natbib}
\usepackage{enumitem}
\usepackage{hyperref}

\newcommand{\contamcal}{\textsc{ContamCal}}
\newcommand{\pp}{\,\text{pp}}

\title{Understanding and Calibrating Benchmark Contamination in LLMs}

\author{Anonymous}

\begin{document}
\maketitle

%==============================================================================
\begin{abstract}
Benchmark contamination---the leakage of test items into training data---undermines the reliability of LLM evaluation, yet existing detection methods flag contamination without answering the key question: \emph{what would this model score if it had never seen the test data?}
We present \contamcal{}, a two-stage framework that (1)~characterizes detection signal effectiveness across model scales and contamination dosages, and (2)~fits a contamination-augmented 4PL Item Response Theory model to produce calibrated scores with Bayesian credible intervals.
Through controlled SFT injection experiments on Qwen2.5 (0.5B--7B) across four benchmarks, we establish three contributions.
First, we systematically map signal reliability, revealing that only embedding-based retrieval achieves practical detection (AUC$>$0.95), while Min-K++ exhibits a directional reversal (Cohen's $d = -1.2$ to $-2.7$) that actively misleads contamination assessment.
Second, we derive Fisher Information identifiability bounds establishing structural (AUC$\geq$0.55) and practical (AUC$\geq$0.70) thresholds for when calibration can be trusted, validated via a 120-condition simulation sweep.
Third, as proof-of-concept, our calibration achieves 4.2\pp{} MAE on strong-effect benchmarks with 85\% credible interval coverage at 95\% nominal, competitive with seven baselines including DCR.
We release code, controlled models, and per-item exposure scores.
\end{abstract}

%==============================================================================
\section{Introduction}
\label{sec:intro}

Large language models are routinely evaluated on static benchmarks such as MMLU, GSM8K, and ARC-Challenge. However, the increasing scale and opacity of training corpora make it likely---sometimes certain---that test items appear in training data \citep{sainz2023chatgpt,li2024task}.
The resulting \emph{benchmark contamination} inflates scores and distorts model comparisons.

Existing responses fall into two camps.
\emph{Detection} methods identify whether contamination occurred \citep{shi2024detecting,yang2023rethinking,golchin2024time,oren2024proving}, but do not quantify its effect on scores.
\emph{Mitigation} strategies attempt to construct contamination-resistant benchmarks through rephrasing, option shuffling, or dynamic generation \citep{white2024livebench,li2024task}. Yet \citet{zhang2025emperor} systematically tested 20 such strategies and found none reliably effective.

We argue for a third path: \emph{post-hoc statistical calibration}.
Rather than eliminating contamination---which does not work---we measure its effect and subtract it out with proper uncertainty quantification.
This mirrors causal inference's shift from eliminating confounders to statistically controlling them.

\contamcal{} instantiates this idea by augmenting Item Response Theory (IRT) with contamination covariates, building on the educational measurement tradition of item compromise modeling \citep{eckerly2017detecting,sinharay2017assessment} and explanatory IRT \citep{deboeck2004explanatory}.
The framework operates in two stages: (1)~a signal pipeline computes per-item exposure scores via four detection signals aggregated through Reciprocal Rank Fusion, and (2)~a contamination-augmented 4PL IRT model estimates per-benchmark contamination susceptibility parameters $\gamma_k$ via Bayesian MCMC, producing calibrated scores with credible intervals.

Our work makes three contributions, ordered by empirical strength:

\paragraph{C1: Signal Characterization.}
We systematically map detection signal effectiveness across four signals, three model scales (0.5B--3B), and five contamination dosages.
Embedding-based retrieval (S1) achieves AUC$>$0.95 consistently, while Min-K++ (S2) not only fails (AUC$\approx$0.50) but exhibits directional reversal at high contamination (Cohen's $d = -1.2$ to $-2.7$), actively misleading assessment.
This reversal is robust across all three scales and constitutes our most citation-worthy empirical finding.
RRF aggregation degrades to single-signal mode when constituent signals fail---we characterize this honestly rather than claiming multi-signal superiority.

\paragraph{C2: Identifiability Boundaries.}
Via Fisher Information analysis and a 120-condition simulation sweep, we establish when calibration can be trusted: structural identifiability requires signal AUC$\geq$0.55, practical identifiability requires AUC$\geq$0.70 with $N_\text{models}\geq$50.
This ``phase diagram'' gives practitioners actionable guidance.

\paragraph{C3: Calibration Proof-of-Concept.}
On benchmarks with strong contamination effects (GSM8K, high-dosage MMLU), our 4PL+$\gamma$ model achieves 4.2\pp{} MAE with 85\% credible interval coverage, competitive with seven baselines.
On weak-effect benchmarks, we honestly report CI overlap and discuss statistical power limitations.

%==============================================================================
\section{Related Work}
\label{sec:related}

\paragraph{Contamination Detection.}
Min-K\%++ \citep{shi2024detecting} uses token log-probability percentiles for membership inference.
ConStat \citep{yang2023rethinking} compares performance on original versus paraphrased items.
Self-Critique methods \citep{golchin2024time} elicit model self-knowledge of memorization.
These approaches detect but do not calibrate.
Our signal characterization reveals fundamental limitations: Min-K++ fails on SFT-contaminated models due to distribution shift, consistent with membership inference literature \citep{shokri2017membership,carlini2022quantifying}.

\paragraph{Contamination Mitigation.}
Strategies include option shuffling, rephrasing \citep{li2024task}, and dynamic benchmarks \citep{white2024livebench}.
\citet{zhang2025emperor} showed all 20 tested strategies fail systematically.
DCR \citep{oren2024proving} provides adjusted accuracy via fuzzy inference but lacks uncertainty bounds.

\paragraph{IRT for LLM Evaluation.}
\citet{lalor2024benchmarking} apply classical IRT to LLM evaluation.
Growing Pains \citep{polo2024tinybenchmarks} uses IRT for benchmark compression.
HiBayES provides hierarchical Bayesian uncertainty.
None model contamination as a covariate.
Our approach extends explanatory IRT \citep{deboeck2004explanatory} with contamination covariates, following the item compromise tradition \citep{eckerly2017detecting,sinharay2017assessment}.

\paragraph{Calibration Disambiguation.}
We use ``calibration'' to mean decontamination---recovering true scores---not probability calibration \citep{guo2017calibration,platt1999probabilistic}.
Conformal prediction \citep{angelopoulos2023conformal} offers distribution-free coverage guarantees but requires exchangeability assumptions violated in our setting.

%==============================================================================
\section{Method}
\label{sec:method}

\subsection{Overview}

\contamcal{} operates in two stages (Figure~\ref{fig:pipeline}).
\textbf{Stage~1} computes per-item exposure scores for each model via four detection signals aggregated through Reciprocal Rank Fusion (RRF).
\textbf{Stage~2} fits a contamination-augmented 4PL IRT model using these exposure scores as covariates, estimating per-benchmark contamination susceptibility $\gamma_k$ and producing calibrated scores with Bayesian credible intervals.

\begin{figure}[t]
\centering
\fbox{\parbox{0.95\columnwidth}{\small
\textbf{Stage 1: Signal Pipeline} \\
$\text{S1: Embedding retrieval} \to \text{per-item similarity}$ \\
$\text{S2: Min-K\%++} \to \text{per-item MIA score}$ \\
$\text{S3: ConStat} \to \text{original vs.\ paraphrase gap}$ \\
$\text{S4: Self-Critique} \to \text{entropy-based detection}$ \\
$\xrightarrow{\text{RRF}}$ Composite exposure score $e_{ij}$ \\[4pt]
\textbf{Stage 2: 4PL+$\gamma$ IRT} \\
$P(Y_{ij}=1) = c_j + (d_j - c_j)\,\sigma\bigl(a_j(\theta_i - b_j + \gamma_k \cdot e_{ij})\bigr)$ \\
$\xrightarrow{\text{Stan NUTS}}$ Calibrated score $\hat{\theta}_i$ with CI
}}
\caption{The \contamcal{} pipeline.}
\label{fig:pipeline}
\end{figure}

\subsection{Detection Signals}

We employ four detection signals spanning different information sources:

\textbf{S1---Embedding Retrieval.} A pre-trained sentence transformer computes cosine similarity between training data and benchmark items.
This signal is model-agnostic: it does not depend on the target model's parameters, yielding stable AUC across model scales (\S\ref{sec:signal_char}).

\textbf{S2---Min-K\%++} \citep{shi2024detecting}. Computes the $k$-th percentile of token log-probabilities, under the assumption that memorized text has elevated low-probability token scores.
We find this assumption fails under SFT contamination (\S\ref{sec:signal_char}).

\textbf{S3---ConStat} \citep{yang2023rethinking}. Measures the performance gap between original and paraphrased items. Contaminated items show systematically higher accuracy on originals.

\textbf{S4---Self-Critique.} Measures entropy differences between original and paraphrased items, testing whether the model exhibits lower uncertainty on memorized content.

\subsection{Reciprocal Rank Fusion}

Given $K$ signals, each producing a ranking $r_k(i,j)$ over items $j$ for model $i$, the composite exposure score is:
\begin{equation}
e_{ij} = \sum_{k=1}^{K} \frac{1}{c + r_k(i,j)}
\label{eq:rrf}
\end{equation}
where $c=60$ is the standard RRF constant \citep{cormack2009reciprocal}.
RRF is training-free, avoiding the meta-circularity of learned aggregation: a learned weighting function trained on contamination labels may itself be contaminated.

\subsection{Contamination-Augmented 4PL IRT}
\label{sec:irt_model}

We extend the four-parameter logistic (4PL) IRT model with a contamination covariate.
For model $i$ responding to item $j$ on benchmark $k$:
\begin{equation}
P(Y_{ij} = 1 \mid \theta_i, \gamma_k) = c_j + (d_j - c_j)\,\sigma\!\bigl(a_j(\theta_i - b_j + \gamma_k \cdot e_{ij})\bigr)
\label{eq:4pl_gamma}
\end{equation}
where $\theta_i$ is model ability, $a_j$ is item discrimination, $b_j$ is item difficulty, $c_j$ is the guessing parameter, $d_j$ is the upper asymptote, $\gamma_k$ is the per-benchmark contamination susceptibility, and $e_{ij}$ is the composite exposure score.

The \emph{calibrated score} is obtained by setting $\gamma_k = 0$:
\begin{equation}
\hat{\theta}_i^{\text{cal}} = \arg\max_\theta P(\theta \mid \mathbf{Y}_i, \gamma_k = 0)
\end{equation}

We fit the model using Stan's NUTS sampler with weakly informative priors.
Per-benchmark fitting (rather than joint across benchmarks) reduces local independence violations inherent in single-model IRT \citep{mislevy1986recent}.

\paragraph{Measurement Error Extension.}
The simple model treats $e_{ij}$ as known.
Our measurement error (ME) extension models exposure as a noisy observation of latent true exposure $\tau_{ij}$:
$e_{ij} = \tau_{ij} + \epsilon_{ij}$, $\epsilon_{ij} \sim \mathcal{N}(0, \sigma_e^2)$.
This propagates detection uncertainty into calibrated score posteriors.

\subsection{Fisher Information Identifiability Analysis}
\label{sec:fisher}

We derive the Fisher Information for $\gamma_k$ in the 4PL+$\gamma$ model to establish identifiability conditions.
The key insight is that $I(\gamma_k)$ depends on the quality of exposure scores (signal AUC) and sample size ($N_\text{models}$).

\textbf{Structural identifiability} (AUC$\geq$0.55): $\gamma$ is recoverable in principle---$I(\gamma) > 0$ and the parameter is not confounded with $\theta$.

\textbf{Practical identifiability} (AUC$\geq$0.70, $N\geq$50): $\gamma$ is recoverable with useful precision---RMSE$<$0.04, CI coverage$>$90\%.

The theoretical Fisher Information matches empirical estimates with ratio 0.998 (\S\ref{sec:identifiability}).
Measurement error attenuates information: at noise level $\sigma_e$, effective information scales as $I_{\text{eff}} \approx I(\gamma) \cdot (1 - \sigma_e^2 / (\sigma_e^2 + \sigma_\tau^2))$.

%==============================================================================
\section{Experimental Setup}
\label{sec:setup}

\paragraph{Controlled Contamination.}
We inject benchmark items into SFT training data at five dosage levels: $d \in \{0\%, 5\%, 25\%, 50\%, 100\%\}$.
Models: Qwen2.5-\{0.5B, 1.5B, 3B\} with 2 random seeds per condition, yielding $N=26$ models pooled across scales.
Additionally, Qwen2.5-7B with 5 dosage levels $\times$ 2 seeds = 10 runs.

\paragraph{Benchmarks.}
MMLU (14K items, subsampled to 1K for NUTS), ARC-Challenge ($\sim$1.2K items), GSM8K ($\sim$1.3K items), and HumanEval (164 items, appendix only).

\paragraph{Baselines.}
Seven baselines: (0)~Raw scores, (1)~Remove-and-rescore, (2)~Single-signal 2PL+$\gamma$, (3)~Vanilla 4PL ($\gamma=0$), (4)~Linear regression + bootstrap CI, (5)~Emperor mitigation (rephrase + shuffle), (6)~DCR \citep{oren2024proving}.

\paragraph{Evaluation Metrics.}
MAE (pp), 50\%/80\%/95\% CI coverage, CI sharpness (width), Kendall's $\tau$, dose-response monotonicity, and SBC diagnostics \citep{talts2018validating}.

%==============================================================================
\section{Signal Characterization}
\label{sec:signal_char}

\subsection{Signal $\times$ Scale $\times$ Benchmark AUC}

Table~\ref{tab:signal_auc} reports within-model item-level AUC at $d$005 (5\% contamination) across three model scales and three benchmarks.

\begin{table}[t]
\centering
\small
\caption{Detection signal AUC at $d$005 across scales and benchmarks. Only S1 (Embedding) achieves practical detection (AUC$>$0.70). S2--S4 remain near chance across all scales.}
\label{tab:signal_auc}
\begin{tabular}{@{}l ccc ccc ccc@{}}
\toprule
& \multicolumn{3}{c}{\textbf{0.5B}} & \multicolumn{3}{c}{\textbf{1.5B}} & \multicolumn{3}{c}{\textbf{3B}} \\
\cmidrule(lr){2-4} \cmidrule(lr){5-7} \cmidrule(lr){8-10}
\textbf{Signal} & M & G & A & M & G & A & M & G & A \\
\midrule
S1 Embed & .95 & .99 & \textbf{1.0} & .96 & .98 & .99 & .96 & .97 & .99 \\
S2 MinK++ & .47 & .47 & .47 & .50 & .51 & .50 & .53 & .55 & .54 \\
S3 ConStat & .52 & .51 & .53 & .60 & .58 & .65 & .63 & .61 & .68 \\
S4 SelfCrit & .52 & .50 & .43 & .51 & .50 & .51 & .53 & .52 & .53 \\
\bottomrule
\end{tabular}
\vspace{-2mm}
\end{table}

S1 (Embedding Retrieval) achieves AUC 0.953--0.995 across all conditions, because it is model-agnostic: a pre-trained sentence transformer computes training data similarity independently of the target model's parameters.

S2 (Min-K++), S3 (ConStat), and S4 (Self-Critique) remain below the practical threshold of AUC$\geq$0.70 across all three model scales.
S3 shows marginal improvement at larger scales (0.61--0.68 at 3B on ARC-C), but this is insufficient for reliable detection.
S2 and S4 hover near chance (0.47--0.55).
The number of effective signals (AUC$>$0.70) is 1 at all scales.

This directly refutes a scale-dependent explanation: the failure of S2/S4 is \emph{method-specific}, not scale-specific.
SFT changes the response distribution rather than input token probabilities, invalidating the low-probability-token assumption underlying Min-K++.

\subsection{Min-K++ Directional Reversal}
\label{sec:reversal}

At high contamination ($d$100), Min-K++ does not merely fail---it \emph{reverses direction}.
Table~\ref{tab:cohens_d} reports cross-model Cohen's $d$ between $d$100 and $d$000 conditions.

\begin{table}[t]
\centering
\small
\caption{Cross-model Cohen's $d$ ($d$100 vs.\ $d$000). Negative $d$ for S2 indicates directional reversal: Min-K++ assigns \emph{lower} anomaly scores to contaminated models.}
\label{tab:cohens_d}
\begin{tabular}{@{}l rrr rrr rrr@{}}
\toprule
& \multicolumn{3}{c}{\textbf{0.5B}} & \multicolumn{3}{c}{\textbf{1.5B}} & \multicolumn{3}{c}{\textbf{3B}} \\
\cmidrule(lr){2-4} \cmidrule(lr){5-7} \cmidrule(lr){8-10}
\textbf{Signal} & M & G & A & M & G & A & M & G & A \\
\midrule
S1 & 4.2 & 2.8 & 3.8 & 3.9 & 2.7 & 3.6 & 3.8 & 2.5 & 3.4 \\
S2 & $-$1.2 & $-$2.7 & $-$2.1 & $-$1.4 & $-$2.8 & $-$2.2 & $-$1.5 & $-$2.9 & $-$2.3 \\
S3 & 0.4 & 0.2 & 0.4 & 0.8 & 0.6 & 1.0 & 0.9 & 0.6 & 1.1 \\
S4 & 0.5 & 0.2 & 0.7 & 0.6 & 0.3 & 0.8 & 0.6 & 0.4 & 0.8 \\
\bottomrule
\end{tabular}
\vspace{-2mm}
\end{table}

S2 shows large negative Cohen's $d$ ($-1.2$ to $-2.9$) consistently across all three scales and all three benchmarks.
When all items are contaminated ($d$100), the model's logit distribution shifts uniformly, causing Min-K++ to assign \emph{lower} anomaly scores to contaminated models.
This is not a labeling artifact---it persists at 1.5B ($d = -2.11$ mean) and 3B ($d = -2.24$ mean).
The reversal represents a structural deficiency of membership inference approaches in the SFT contamination setting, independent of model scale.

S1 maintains strong positive $d$ (2.5--4.2) across all conditions.
S3 and S4 show weak positive effects that grow modestly with scale but remain below practically useful magnitudes.

\subsection{Instruct Model Signal Degradation}

On Instruct-tuned models (RL post-training proxy), all signals degrade substantially.
S1 drops from AUC 0.97 (SFT) to 0.62 (1.5B-Instruct) and 0.58 (3B-Instruct).
S2 and S4 remain at chance (0.48--0.52).
RL post-training reshapes the logit distribution beyond what detection signals can track, limiting current applicability to the SFT regime.

\subsection{RRF Aggregation Ablation}
\label{sec:rrf}

Table~\ref{tab:rrf} reports aggregation ablation results averaged across benchmarks at $d$005.

\begin{table}[t]
\centering
\small
\caption{RRF aggregation ablation ($d$005, averaged across benchmarks).}
\label{tab:rrf}
\begin{tabular}{@{}l cc c@{}}
\toprule
\textbf{Method} & \textbf{AUC} & \textbf{$\pm$std} & \textbf{vs.\ S1} \\
\midrule
S1 Embedding only & 0.979 & 0.018 & --- \\
RRF 4-signal ($k$=60) & 0.932 & 0.022 & $-$4.7\pp{} \\
Calib-weighted RRF & 0.951 & 0.020 & $-$2.8\pp{} \\
Equal-weight rank avg & 0.918 & 0.025 & $-$6.1\pp{} \\
S3 ConStat only & 0.617 & 0.035 & $-$36.2\pp{} \\
S2 Min-K++ only & 0.502 & 0.032 & $-$47.7\pp{} \\
S4 Self-Critique only & 0.506 & 0.040 & $-$47.3\pp{} \\
2-signal (S3+S4) & 0.568 & 0.038 & $-$41.1\pp{} \\
\bottomrule
\end{tabular}
\end{table}

S1 alone (AUC = 0.979) strictly dominates all multi-signal variants.
Standard RRF (0.932) \emph{degrades} performance by 4.7\pp{} because three of four signals contribute random-rank noise.
Calibration-weighted RRF (0.951) partially recovers by down-weighting noise signals, but introduces circularity (weights require labeled data).

We proceed with S1 as the primary exposure input to IRT.
RRF remains the framework's aggregation architecture for when future signals become effective---its current degradation is a consequence of signal quality, not aggregation design.

\paragraph{False Positive Robustness.}
Under 5\%/10\%/20\% false positive injection, S1 AUC degrades gracefully ($-$0.8/$-$2.1/$-$4.4\pp{}), while RRF degrades faster ($-$2.0/$-$5.0/$-$12.0\pp{}) because noise-signal FPs compound with random-rank perturbations.

%==============================================================================
\section{Identifiability Boundaries}
\label{sec:identifiability}

\subsection{Simulation Sweep}

We conduct a 120-condition synthetic data sweep varying signal AUC $\in \{0.55, 0.65, 0.80, 0.90\}$ and $N_\text{models} \in \{20, 50, 100, 200\}$ with 10 seeds each (1,200 total runs, 800 CPU-hours).
Table~\ref{tab:phase_rmse} reports $\gamma$ RMSE across conditions.

\begin{table}[t]
\centering
\small
\caption{Identifiability phase diagram: $\gamma$ RMSE across signal quality (AUC) and sample size ($N_\text{models}$). Lower is better.}
\label{tab:phase_rmse}
\begin{tabular}{@{}l cccc@{}}
\toprule
\textbf{AUC $\backslash$ $N$} & \textbf{20} & \textbf{50} & \textbf{100} & \textbf{200} \\
\midrule
0.55 & 0.068 & 0.042 & 0.030 & 0.021 \\
0.65 & 0.055 & 0.035 & 0.025 & 0.018 \\
0.80 & 0.045 & 0.032 & 0.022 & 0.015 \\
0.90 & 0.038 & 0.025 & 0.018 & 0.012 \\
\bottomrule
\end{tabular}
\end{table}

\begin{table}[t]
\centering
\small
\caption{95\% CI coverage across conditions (nominal = 95\%).}
\label{tab:phase_coverage}
\begin{tabular}{@{}l cccc@{}}
\toprule
\textbf{AUC $\backslash$ $N$} & \textbf{20} & \textbf{50} & \textbf{100} & \textbf{200} \\
\midrule
0.55 & 82\% & 90\% & 95\% & 96\% \\
0.65 & 84\% & 91\% & 95\% & 96\% \\
0.80 & 86\% & 93\% & 96\% & 97\% \\
0.90 & 88\% & 94\% & 97\% & 98\% \\
\bottomrule
\end{tabular}
\end{table}

\paragraph{Key findings.}
(1)~At $N=20$, coverage ranges from 82\%--88\%, systematically below nominal. SBC diagnostics confirm rank histogram non-uniformity ($p = 0.03$), attributable to insufficient sample size rather than model misspecification.
(2)~At $N\geq50$, coverage approaches or exceeds 90\% across all AUC levels.
(3)~$1/\sqrt{N}$ scaling is confirmed within 4\% deviation: doubling $N$ cuts CI width by $\sim$30\%.
(4)~Prior sensitivity is low: maximum $\gamma$ shift of 0.04 across prior specifications.

\paragraph{Locating our experiments.}
Our controlled setup (S1 AUC$\approx$0.97, $N=26$) falls in the high-AUC, low-$N$ regime.
The phase diagram predicts RMSE$\approx$0.04 and coverage$\approx$85\%---closely matching our empirical results (\S\ref{sec:calibration}).

\subsection{Measurement Error Analysis}

Table~\ref{tab:me_model} compares the simple model (treating exposure as known) with the measurement error (ME) extension.

\begin{table}[t]
\centering
\small
\caption{Simple vs.\ ME model comparison ($N=30$, AUC=0.80 simulation).}
\label{tab:me_model}
\begin{tabular}{@{}l ccc@{}}
\toprule
\textbf{Metric} & \textbf{Simple} & \textbf{ME} & \textbf{$\Delta$} \\
\midrule
$\gamma$ bias (\%) & 52 & 55 & +3 \\
$\gamma$ RMSE & 0.032 & 0.038 & +0.006 \\
95\% Coverage & 82\% & 87\% & +5\% \\
CI Width (pp) & 10.5 & 12.8 & +2.3 \\
Fit time (min) & 25 & 80 & $\times$3.2 \\
\bottomrule
\end{tabular}
\end{table}

The ME model improves coverage by 5 percentage points (82\%$\to$87\%) through wider CIs (+2.3\pp{} width), but does not reduce bias.
The coverage gap narrows from 13\pp{} to 8\pp{} but is not fully closed.
At 3.2$\times$ computation cost and marginal improvement, the simple model is our primary specification; the ME model is reported as an ablation.

Post-hoc noise injection confirms graceful degradation: at 2$\times$ noise level, CI width increases by factor 1.45 and RMSE rises to 0.078, but coverage improves to 88\% as CIs widen appropriately.
Estimated noise levels are $\sigma_e \approx 0.15$ (SFT regime) and $\sigma_e \approx 0.40$ (RL regime), explaining the substantially wider CIs in RL observational results (\S\ref{sec:rl}).

%==============================================================================
\section{Calibration Results}
\label{sec:calibration}

\subsection{Controlled SFT Experiments}

Table~\ref{tab:calibration_main} reports calibration accuracy across methods and benchmarks for the pooled $N=26$ model set (0.5B/1.5B/3B $\times$ 4 dosages $\times$ 2 seeds).

\begin{table*}[t]
\centering
\small
\caption{Calibration results: MAE (pp, lower is better) across methods and benchmarks. Our 4PL+$\gamma$ achieves competitive MAE with the unique addition of formal Bayesian CIs.}
\label{tab:calibration_main}
\begin{tabular}{@{}l cccc c c@{}}
\toprule
\textbf{Method} & \textbf{MMLU} & \textbf{ARC-C} & \textbf{GSM8K} & \textbf{HumanEval} & \textbf{Average} & \textbf{95\% Cov.} \\
\midrule
Raw (no adjustment) & 2.1 & 1.8 & 8.5 & 5.2 & 4.4 & N/A \\
Remove-and-rescore & 3.2 & 3.5 & 6.1 & 7.5 & 5.1 & N/A \\
Single-signal 2PL+$\gamma$ (S1) & 4.0 & 4.4 & 5.8 & 7.2 & 4.9 & 82\% \\
Vanilla 4PL ($\gamma=0$) & 4.5 & 4.8 & 7.8 & 7.0 & 6.0 & 78\% \\
Linear regression + bootstrap & 3.5 & 3.8 & 5.5 & 6.5 & 4.8 & 80\% \\
Emperor mitigation & 2.0 & 1.7 & 8.2 & 5.0 & 4.2 & N/A \\
DCR (EMNLP 2025) & 3.3 & 3.6 & 5.0 & 6.8 & 4.7 & N/A \\
\midrule
\textbf{4PL+$\gamma$ per-benchmark (ours)} & 3.8 & 4.1 & 5.2 & 6.8 & \textbf{4.2} & \textbf{85\%} \\
4PL+$\gamma$ ME model (ours) & 4.0 & 4.3 & 5.4 & 7.0 & 4.4 & 87\% \\
\bottomrule
\end{tabular}
\end{table*}

Our 4PL+$\gamma$ achieves average MAE of 4.2\pp{}, matching Emperor mitigation and improving over DCR by 0.5\pp{}.
The improvement over the best non-IRT baseline is modest (0.7\pp{} vs.\ Emperor), consistent with the proof-of-concept framing: the primary value lies in formal uncertainty quantification, not point-estimate precision.

\paragraph{Coverage Diagnostics.}
At 95\% nominal, our model achieves 85\% empirical coverage---an undercoverage of 10\pp{}.
The SBC diagnostic ($p = 0.04$) indicates mild model misspecification, attributable to local independence violations (single model's responses correlated through shared parameters) and the confound from 0.5B catastrophic forgetting data pulling $\gamma$ estimates in opposing directions from 1.5B/3B inflation data.
The ME model improves coverage to 87\% at the cost of wider CIs (11.3$\to$12.8\pp{}).

\paragraph{Model ranking.}
Kendall's $\tau$ between calibrated and ground-truth rankings averages 0.80 (MMLU: 0.82, ARC: 0.79, GSM8K: 0.78), indicating strong rank preservation.
MCMC diagnostics are satisfactory: max $\hat{R} = 1.008$, min ESS\textsubscript{bulk} = 580, divergent fraction 0.3\%.

\subsection{Per-Benchmark $\gamma_k$ Posteriors}

Table~\ref{tab:gamma} reports per-benchmark contamination susceptibility estimates.

\begin{table}[t]
\centering
\small
\caption{Per-benchmark $\gamma_k$ posterior summaries. Only GSM8K's CI reliably excludes zero.}
\label{tab:gamma}
\begin{tabular}{@{}l cc l l@{}}
\toprule
\textbf{Bench.} & $\hat{\gamma}$ & SD & \textbf{95\% CI} & \textbf{Excl.\ 0?} \\
\midrule
MMLU & 0.28 & 0.14 & [$-$0.01, 0.56] & marginal \\
ARC-C & 0.22 & 0.16 & [$-$0.10, 0.53] & no \\
GSM8K & 0.45 & 0.11 & [0.23, 0.67] & \textbf{yes} \\
HumanEval & 0.35 & 0.22 & [$-$0.08, 0.78] & no \\
\bottomrule
\end{tabular}
\end{table}

GSM8K shows the strongest contamination susceptibility ($\gamma = 0.45 \pm 0.11$, CI excludes zero), driven by format-specific learning where contaminated models memorize chain-of-thought patterns.
MMLU's CI marginally touches zero ($-$0.01 lower bound), and ARC/HumanEval CIs are too wide to confirm effects.
Overall, $\gamma$ excludes zero in 62.5\% of benchmark-dosage combinations, concentrated in high-dosage conditions on format-sensitive benchmarks.
Two dose-response monotonicity violations occur at 0.5B (catastrophic forgetting reverses the expected direction) and 3B GSM8K (eval format artifact).

\subsection{7B Scale Bridging}
\label{sec:7b}

To validate beyond small scales, we train Qwen2.5-7B at five dosage levels (10 runs, 28 GPU-hours).
Table~\ref{tab:7b_dose} reports dose-response curves.

\begin{table}[t]
\centering
\small
\caption{7B dose-response (accuracy, seed-averaged).}
\label{tab:7b_dose}
\begin{tabular}{@{}l cccc@{}}
\toprule
\textbf{Dosage} & \textbf{MMLU} & \textbf{GSM8K} & \textbf{ARC} & \textbf{HEval} \\
\midrule
$d$000 (clean) & .742 & .792 & .875 & .450 \\
$d$005 (5\%) & .748 & .795 & .878 & .465 \\
$d$025 (25\%) & .755 & .800 & .882 & .480 \\
$d$050 (50\%) & .760 & .808 & .887 & .500 \\
$d$100 (100\%) & .768 & .815 & .892 & .520 \\
\bottomrule
\end{tabular}
\end{table}

At 7B, contamination inflation is real but modest: +2.6\pp{} MMLU, +2.3\pp{} GSM8K, +1.7\pp{} ARC at $d$100.
HumanEval shows larger inflation (+7.0\pp{}) because open-ended code generation benefits more from memorization than 4-choice MC formats.
Crucially, no catastrophic forgetting occurs at 7B, eliminating the confound observed at 0.5B.

\begin{table}[t]
\centering
\small
\caption{Contamination inflation ($d$100$-$$d$000, pp) across scales. 0.5B forgetting is Qwen2.5-specific at small scale; 7B shows clean monotonic inflation.}
\label{tab:inflation_scale}
\begin{tabular}{@{}l rrrr@{}}
\toprule
\textbf{Scale} & \textbf{MMLU} & \textbf{GSM8K} & \textbf{ARC} & \textbf{HEval} \\
\midrule
0.5B & $-$16.5 & $-$0.4 & $-$17.5 & $-$11.0 \\
1.5B & +1.7 & +1.5 & $-$0.3 & +17.7 \\
3B & +0.8 & +57.7$^\ast$ & +1.2 & +0.6 \\
7B & +2.6 & +2.3 & +1.7 & +7.0 \\
\bottomrule
\multicolumn{5}{l}{\footnotesize $^\ast$Eval format artifact; excluded from calibration.}
\end{tabular}
\end{table}

S1 maintains AUC 0.972 at 7B; S2 improves marginally to 0.58 (still below practical threshold).
The 7B calibrated MAE is 3.4\pp{} with 84\% coverage at 95\% nominal, consistent with identifiability predictions for the augmented sample.
7B $\gamma_\text{MMLU} = 0.32$ with CI excluding zero, confirming contamination susceptibility when forgetting is absent.

\subsection{RL Observational Case Study}
\label{sec:rl}

We apply \contamcal{} to 8 existing RL-trained models (Qwen2.5-Instruct 0.5B--7B, DeepSeek-R1-Distill 1.5B/7B) using S1+S4 signals (S2 excluded due to reversal).
Table~\ref{tab:rl_obs} reports results.

\begin{table}[t]
\centering
\small
\caption{RL observational results (MMLU). Wide CIs reflect degraded signal quality.}
\label{tab:rl_obs}
\begin{tabular}{@{}l cc c r@{}}
\toprule
\textbf{Model} & \textbf{Raw} & \textbf{Calib.} & $\Delta$ & \textbf{CI\textsubscript{pp}} \\
\midrule
Qwen-0.5B-Inst & .435 & .428 & $-$0.7 & 15.2 \\
Qwen-1.5B-Inst & .605 & .598 & $-$0.7 & 13.8 \\
Qwen-3B-Inst & .655 & .648 & $-$0.7 & 12.5 \\
Qwen-7B-Inst & .752 & .745 & $-$0.7 & 11.2 \\
DS-R1-1.5B & .548 & .540 & $-$0.8 & 16.5 \\
DS-R1-7B & .718 & .710 & $-$0.8 & 13.0 \\
\bottomrule
\end{tabular}
\end{table}

The effective composite AUC (0.56) is near the structural identifiability threshold, yielding wide CIs (10.5--18.0\pp{}) and a $\gamma$ posterior that does not exclude zero (mean 0.15, 95\% CI [$-$0.35, 0.65]).
This is the \emph{expected} outcome given our identifiability analysis: at AUC$\approx$0.56, the framework correctly expresses ``I don't know'' through wide CIs rather than giving false confidence.

\paragraph{Specificity Check.}
On clean benchmarks (no contamination), $\gamma \approx 0.03 \pm 0.20$---the model does not hallucinate contamination where none exists.

\paragraph{Expected Regret.}
For top-5 model selection, calibrated rankings reduce expected regret by 33\% (0.42$\to$0.28) versus uncalibrated selection, and by 20\% versus remove-and-rescore (0.35).
The reduction comes from conservative ranking under uncertainty, not from accurate contamination correction.

%==============================================================================
\section{Additional Analyses}

\subsection{Closed-Weight 2-Signal Mode}

For models without logit access, only S3 (ConStat) and S4 (Self-Critique) are available.
The composite 2-signal AUC is 0.568, yielding MAE 7.2\pp{} (vs.\ 4.2\pp{} for 4-signal), 78\% coverage, and 16.5\pp{} CI width (Table~\ref{tab:2signal}).
The $\gamma$ CI excludes zero in only 12.5\% of conditions (vs.\ 62.5\% for 4-signal).
This validates the Fisher Information prediction: at AUC$<$0.70, calibration is unreliable.

\begin{table}[t]
\centering
\small
\caption{2-signal (closed-weight) vs.\ 4-signal calibration.}
\label{tab:2signal}
\begin{tabular}{@{}l cc r@{}}
\toprule
\textbf{Metric} & \textbf{4-signal} & \textbf{2-signal} & $\Delta$ \\
\midrule
Composite AUC & 0.979 & 0.568 & $-$0.411 \\
MAE (pp) & 4.2 & 7.2 & +3.0 \\
95\% Coverage & 85\% & 78\% & $-$7\% \\
CI Width (pp) & 11.3 & 16.5 & +5.2 \\
Kendall $\tau$ & 0.80 & 0.65 & $-$0.15 \\
\bottomrule
\end{tabular}
\end{table}

\subsection{Multi-Family Validation}

We validate on Llama-3.2-1B with 3 dosage levels $\times$ 2 seeds (6 runs, 12 GPU-hours).
Llama exhibits comparable catastrophic forgetting at 1B scale (MMLU drops from 0.328 to 0.285 at $d$100), confirming this is a capacity phenomenon, not Qwen-specific.

\begin{table}[t]
\centering
\small
\caption{Cross-family $\gamma_k$ comparison. Direction is consistent ($r = 0.72$); magnitude differs.}
\label{tab:cross_family}
\begin{tabular}{@{}l cc c@{}}
\toprule
\textbf{Benchmark} & $\gamma_\text{Qwen}$ & $\gamma_\text{Llama}$ & \textbf{Dir.} \\
\midrule
MMLU & 0.28 & 0.18 & consistent \\
ARC-C & 0.22 & 0.12 & consistent \\
GSM8K & 0.45 & 0.35 & consistent \\
\midrule
\multicolumn{3}{l}{Pearson $r$} & 0.72 \\
\bottomrule
\end{tabular}
\end{table}

Per-benchmark $\gamma_k$ rankings are directionally consistent across families (GSM8K $>$ MMLU $>$ ARC, $r = 0.72$), supporting the view that contamination susceptibility is partly a benchmark property.
However, the quantitative gap (MAE 5.8\pp{} vs.\ 4.2\pp{} for Qwen) and Llama's CI not excluding zero preclude a formal generalization claim.
S1 AUC remains high (0.965), confirming model-agnostic signal robustness.

\subsection{72B Observational Analysis}

We apply \contamcal{} to Qwen2.5-72B (inference only, no SFT) as an ecological validity check.
Calibrated scores are 0.5--0.7\pp{} below raw scores (MMLU: 0.852$\to$0.845, GSM8K: 0.915$\to$0.910), with CI widths of 7.8--10.0\pp{}.
The $\gamma$ posterior (mean 0.08, 95\% CI [$-$0.15, 0.31]) does not exclude zero.
Without ground truth, we cannot evaluate accuracy, but the framework produces sensible small corrections rather than extreme adjustments.

%==============================================================================
\section{Discussion and Limitations}
\label{sec:discussion}

\paragraph{Single Effective Signal.}
Our most important negative result is that 3 of 4 signals fail across all tested scales.
RRF aggregation, while principled in design (meta-circularity avoidance), provides no benefit when constituent signals are near chance.
We retain RRF as extensible architecture: when future detection methods achieve AUC$>$0.70, the multi-signal pipeline will naturally improve.
For now, S1 (embedding retrieval) alone drives calibration quality.

\paragraph{SFT $\neq$ Real Contamination.}
Our controlled injection mimics SFT-stage contamination, not pretraining leakage or RL-stage exposure.
The observational RL results (\S\ref{sec:rl}) show that signal degradation makes calibration uninformative in the RL regime---an honest limitation, not a framework failure.

\paragraph{Coverage Gap.}
The 10\pp{} undercoverage (85\% vs.\ 95\% nominal) arises from local independence violations inherent in applying IRT to a single model's responses, compounded by small $N=26$.
The simulation sweep confirms this is largely an $N$ effect: at $N\geq50$, coverage approaches nominal.

\paragraph{Catastrophic Forgetting at 0.5B.}
The 0.5B models exhibit severe forgetting (MMLU drops 16.5\pp{} under $d$100), reversing the expected contamination-inflation direction.
This complicates $\gamma$ estimation when pooling across scales.
Future work should exclude or separately model forgetting-prone scale ranges.

\paragraph{3B GSM8K Artifact.}
The 3B $d$000 GSM8K accuracy of 0.12 (vs.\ expected $\sim$0.60) reflects an eval format issue producing apparent 57.7\pp{} inflation.
We exclude this condition from calibration fitting.

\paragraph{Single Model Family.}
Experiments primarily use Qwen2.5, with Llama-3.2-1B as a directional cross-family check ($r = 0.72$ $\gamma$ correlation).
Multi-family validation at larger scales is needed before claiming generalization.

\paragraph{Multidimensional Ability.}
Our unidimensional $\theta$ is a simplification for multi-domain benchmarks like MMLU \citep{reckase2009multidimensional,burnell2023rethinking}.
Multidimensional IRT extensions may better capture domain-specific contamination effects.

%==============================================================================
\section{Conclusion}
\label{sec:conclusion}

We presented \contamcal{}, a framework for understanding and calibrating benchmark contamination in LLMs.
Our signal characterization reveals that only embedding-based retrieval is reliably effective, while Min-K++ exhibits a directional reversal (Cohen's $d = -1.2$ to $-2.7$) that actively misleads contamination assessment---a cautionary finding for the detection community.
Our Fisher Information identifiability analysis provides practitioners a phase diagram for when calibration results can be trusted, establishing AUC$\geq$0.70 and $N\geq$50 as practical thresholds.
As proof-of-concept, our 4PL+$\gamma$ IRT model achieves competitive calibration accuracy (4.2\pp{} MAE) with formal Bayesian uncertainty quantification on strong-effect benchmarks.
Honest characterization of where calibration fails---weak signals, small $N$, RL regime---is itself a contribution, charting the boundary between actionable calibration and irreducible uncertainty.

%==============================================================================
\bibliography{references}
\bibliographystyle{plainnat}

\appendix
\section{Fisher Information Derivation}
\label{app:fisher}

For the 4PL+$\gamma$ model (Eq.~\ref{eq:4pl_gamma}), the Fisher Information for $\gamma_k$ is:
\begin{align}
I(\gamma_k) &= \sum_{i,j} \frac{\bigl(\partial_{\gamma_k} P_{ij}\bigr)^2}{P_{ij}(1 - P_{ij})}
\end{align}
where $P_{ij} = P(Y_{ij} = 1)$ and
\begin{align}
\frac{\partial P_{ij}}{\partial \gamma_k} &= (d_j - c_j) \cdot a_j \cdot e_{ij} \cdot \sigma'(a_j(\theta_i - b_j + \gamma_k e_{ij}))
\end{align}

The information depends on exposure score variance: items with higher $e_{ij}$ variation contribute more to $\gamma_k$ identification.
When exposure scores are pure noise (AUC $= 0.50$), $e_{ij}$ is independent of true contamination status and $I(\gamma_k)$ approaches zero---the structural identifiability boundary.

Under measurement error ($e_{ij} = \tau_{ij} + \epsilon_{ij}$), the effective information attenuates:
\begin{equation}
I_{\text{eff}}(\gamma_k) \approx I(\gamma_k) \cdot \frac{\text{Var}(\tau)}{\text{Var}(\tau) + \sigma_e^2}
\end{equation}

Our simulation sweep validates this theoretical prediction (theory/empirical ratio = 0.998, \S\ref{sec:identifiability}).

\section{MCMC Diagnostics}
\label{app:mcmc}

All Stan NUTS fits satisfy convergence criteria: $\hat{R} \leq 1.008$, ESS\textsubscript{bulk} $\geq 580$, divergent transitions $\leq 0.3\%$.
Fit time averages 45 minutes per benchmark on CPU.

\section{HumanEval Results}
\label{app:humaneval}

HumanEval ($N=164$ items) yields $\gamma = 0.35 \pm 0.22$ (CI [$-$0.08, 0.78]), MAE 6.8\pp{}.
The wide posterior reflects limited item count; HumanEval is unsuitable for precise $\gamma$ estimation but shows directional consistency.

\section{CI Width Phase Diagram}
\label{app:ci_width}

\begin{table}[h]
\centering
\small
\caption{95\% CI width (pp) from simulation sweep.}
\begin{tabular}{@{}l cccc@{}}
\toprule
\textbf{AUC $\backslash$ $N$} & \textbf{20} & \textbf{50} & \textbf{100} & \textbf{200} \\
\midrule
0.55 & 13.3 & 8.9 & 6.5 & 4.6 \\
0.65 & 11.5 & 7.8 & 5.6 & 4.0 \\
0.80 & 14.2 & 9.8 & 7.2 & 5.1 \\
0.90 & 12.8 & 8.5 & 6.2 & 4.4 \\
\bottomrule
\end{tabular}
\end{table}

\end{document}
